% 第六章 总结与展望
\chapter{总结与展望}

\section{总结}

本文通过微磁学仿真，系统分析了阻挫交换铁磁体系中霍普夫子的稳定条件，研究了其自旋波驱动下的动力学行为。

在初始态构建方面，利用局域旋转矩阵与环面坐标映射，实现了任意拓扑荷及反铁磁背景下霍普夫子的理想初态构型的生成。并且在同一套仿真参数条件下，结果表明，无论初始尺寸如何，霍普夫子最终都会弛豫至固定的平衡尺寸（$R_{\text{eq}} = \SI{2.60}{nm}$，$r_{\text{eq}} = \SI{2.16}{nm}$），且这一过程不受背景磁化方向的影响。通过对各向异性参数的测试，发现在本文研究的参数背景下，维持霍普夫子三维纽结结构的临界阈值在 $\num{52e3}$ 至 $\num{55e3} \, \si{J/m^3}$ 之间。

在动力学特性方面，自旋波驱动存在明显的方向选择性。只有面内极化的自旋波能有效推动霍普夫子，面外极化几乎不起作用。不同传播方向的自旋波会引起不同的运动轨迹。面内传播在 $\SI{1000}{GHz}$ 附近引发最大的横向位移，而轴向传播在 $\SI{1100}{GHz}$ 处引起强烈的共振并导致纵向位移。利用这种方向和频率上的差异，只需切换单个波源的驱动频率（如 $\SI{100}{GHz}$ 与 $\SI{1100}{GHz}$ 切换），就能实现霍普夫子沿 $z$ 轴的双向往返运动。这些表现出的频率选择、方向响应等特性，与生物神经元的生物处理过程高度相似，为设计基于自旋波的神经形态器件提供了直接的理论参考。


\section{展望}

根据当前的研究结果，后续工作可向多个物理维度与应用层面进一步拓展。

在神经形态器件的实现探索上，本研究提出的概念器件框架侧重于理论设想，还缺少具体的仿真实现。后续需要重点研究怎么样能够通过仿真实现单个的人工神经元的搭建以及串联，如何在室温环境下生成能支撑霍普夫子稳定存在的材料体系，最后在具体器件集成层面也需要考虑到工业界的实际执行性和可行性。

对于高阶与混合拓扑构型的动力学特征，除基础拓扑荷（$Q_H = 1$）外，高阶 Hopf 指数构型、分数霍普夫子的动力学行为仍有待研究。可以通过已有的第三章研究内容，生成多种不同拓扑荷构型的霍普夫子，作为研究基础，未来的工作可以研究这种复杂形态在自旋波作用下的频率选择性及其它动力学特性。

反铁磁（AFM）背景下 霍普夫子 的动力学行为也有着丰富的研究前景。第三章所建立的初始化方法已经实现生成任意反铁磁背景情况的霍普夫子，相应的仿真初态已经具备基础，并且可以借鉴已有研究进行进一步拓展\cite{zhang2023magnon,pershoguba2021electronic}。反铁磁情况在高速响应与低功耗驱动方面有着不同的潜力。此类体系中 霍普夫子 对自旋波激励的平动规律、方向选择性以及双向控制机制的具体形态，都可以借助系统扫描实现。最后可以将将上述结果与本文铁磁体系所得规律放在一起比较，有助于分析磁化背景对 霍普夫子 动力学特征所施加的作用规律。

此外，多物理场同时考虑与多参数空间的交叉影响也是完善理论框架的重要环节。后续仿真可扩展参数扫描范围，系统考察高阶阻挫交换系数、考虑电场磁场电流驱动以及其他情况的几何边界对拓扑结构稳定性的协同作用。探明上述次级参数的交叉机制，将为相关自旋电子器件的实验制备提供更为完整的参考依据。
